\documentclass[11pt,letterpaper]{article}
\usepackage[margin=0.82in,headheight=15pt]{geometry}
\usepackage{fontspec}
\setmainfont{Latin Modern Roman}
\setsansfont{Latin Modern Sans}
\setmonofont{DejaVu Sans Mono}[Scale=0.82]
\usepackage{amsmath,amssymb,amsthm,mathtools}
\usepackage{booktabs,longtable,array,enumitem,microtype}
\usepackage[dvipsnames]{xcolor}
\usepackage{fvextra,fancyhdr,needspace,xurl,hyperref}
\definecolor{ink}{HTML}{16324F}
\definecolor{accent}{HTML}{16656B}
\hypersetup{colorlinks=true,linkcolor=ink,urlcolor=accent,citecolor=accent,
  pdftitle={Shorter Lean Proofs for Diophantine Consistency},
  pdfauthor={OpenAI, prepared for Vladimir Reshetnikov}}
\pagestyle{fancy}
\fancyhf{}
\fancyhead[L]{\small\sffamily Diophantine Consistency / Lean proof review}
\fancyhead[R]{\small\sffamily 26 October 2026}
\fancyfoot[C]{\small\thepage}
\setlength{\parindent}{0pt}
\setlength{\parskip}{5pt}
\setlist{itemsep=2pt,topsep=4pt,leftmargin=*}
\renewcommand{\arraystretch}{1.16}
\setcounter{tocdepth}{1}
\raggedbottom
\newcommand{\code}[1]{\nolinkurl{#1}}
\newcommand{\pin}{60e9c5d620c6c387d932876cfa2d9376f64ea56c}
\newcommand{\src}[2]{\href{https://github.com/VladimirReshetnikov/ProveIt/blob/60e9c5d620c6c387d932876cfa2d9376f64ea56c/Computability/DiophantineConsistency/Lean/DiophantineConsistency/#1.lean}{\code{#1.lean}}#2}
\DefineVerbatimEnvironment{Lean}{Verbatim}{breaklines=true,breakanywhere=true,samepage=true,
  fontsize=\small,frame=single,rulecolor=\color{black!18},framesep=6pt}
\newtheorem{proposition}{Proposition}

\begin{document}
\begin{titlepage}
\vspace*{0.7in}
{\color{accent}\sffamily\large SOURCE-GROUNDED REFACTORING STUDY\par}
\vspace{0.3in}
{\color{ink}\sffamily\bfseries\Huge Shorter Lean Proofs for\\[4pt] Diophantine Consistency\par}
\vspace{0.3in}
{\Large Concrete simplifications, reusable interfaces,\\checked prototypes, and a migration plan\par}
\vspace{0.45in}
Prepared for Vladimir Reshetnikov\\
26 October 2026
\vspace{0.3in}

\textbf{Reviewed repository.}
\href{https://github.com/VladimirReshetnikov/ProveIt}{VladimirReshetnikov/ProveIt},
\code{Computability/DiophantineConsistency}.\\
\textbf{Pinned commit.} \code{60e9c5d620c6c387d932876cfa2d9376f64ea56c}.\\
\textbf{Toolchain.} Lean 4.32.0; Mathlib revision
\code{81a5d257c8e410db227a6665ed08f64fea08e997}.

\vspace{0.4in}
\textbf{Main recommendation.}
Start with existing library lemmas, common arithmetic identities, and the
interfaces between proof stages. Then extract the common circuit tail and
scanner simulation. These changes remove repeated reasoning while preserving
the exact executable syntax, finite certificates, and theorem hypotheses.

The package contains reproducible source measurements, concrete Lean examples,
patches, and a record distinguishing successful kernel checks from proposed
integrations. It does not claim a completed refactoring of the whole repository
or a new end-to-end consistency theorem.
\vfill
{\small This report reviews the active Lean tree, not every archived manuscript
copy. All quantitative observations refer to the pinned snapshot.}
\end{titlepage}

\section*{Executive assessment}
\addcontentsline{toc}{section}{Executive assessment}

The best opportunities are concentrated in repeated mathematical interfaces:
encoded-list scanners, coordinate renamings, finite coefficient rows, and the
composition of already certified circuit blocks. A uniform substitution of
\code{simp}, \code{aesop}, or a larger heartbeat limit would miss these targets.

There are useful changes at three scales. First, seven private list-lookup
proofs can use two existing Lean-core lemmas; the supplied patch removes
40 physical lines. Second, a common arithmetic identity removes the
half-radix case analysis, and one row-wise argument can replace hundreds of
lines of repeated coefficient bounds. Third, the Split and Gap natural
circuits have an exactly matching 99-node suffix under a specified register
renaming. A reusable block for that suffix is the most concrete larger
architectural target identified here.

Some abstractions are worthwhile without shortening the file immediately.
For example, a common topological-DAG invariant theorem can replace repeated
strong inductions. After charging for the new helper, the first few clients
do not yet yield a net line saving. The report states this explicitly instead
of treating every abstraction as a compression result.

Several elementary examples were checked with the pinned Lean compiler.
Larger replacements and full scanner/circuit migrations have separate status
labels in Section~\ref{sec:validation} and the archive's verification record.
The strongest source comparisons are accompanied by scripts and exact data.

\newpage
\tableofcontents
\newpage

\section{Scope, baseline, and what should be preserved}
\label{sec:scope}

The inspected snapshot is the main-branch commit recorded above. The active
source root is \code{Computability/DiophantineConsistency/Lean}. Its physical
source counts are:

\begin{center}
\begin{tabular}{lrr}
\toprule
Category & Files & Physical lines\\
\midrule
Library modules and aggregate import file & 326 & 66,857\\
Checks and audits & 220 & 7,114\\
Entire active Lean tree & 546 & 73,971\\
\bottomrule
\end{tabular}
\end{center}

The classification puts files under \code{checks/}, and files whose names end
in \code{Checks} or \code{Audit}, in the second row. Physical lines include
comments and blank lines. These are source-volume figures, not counts of
independent mathematical arguments or kernel proof-term sizes. The supplied
\code{scripts/audit_sources.py} reproduces the census, token counts, hashes,
import edges, and approximate declaration spans.

The two largest generated step files account for 6,567 lines:
\src{SplitCircuitSteps59}{ has 3,445, and }
\src{GapCircuitSteps66}{ has 3,122.}
They specialize existing generic evaluation theorems. Their size is relevant
to emitted code, build granularity, and review ergonomics, but should not be
reported as thousands of independently maintained proofs.

After comments and strings are masked, the active source census finds no
\code{sorry}, \code{axiom}, or \code{native_decide} tokens. This lexical result
is not an audit of imported axioms. Endpoint axiom checks remain necessary.

\subsection{Preserve the existing mathematical contract}

The current formalization plan distinguishes older checked endpoints from
the newer literal polynomial, native decoding, and internal-ZFC composition
work. The README and \code{LEAN_EQUIVALENCE_PLAN.md} explicitly retain open
links. A refactoring must preserve these distinctions; shortening a conditional
lemma is not completion of the whole consistency equivalence.\cite{readme,plan}

In particular, preserve the witness domain, the literal coefficient data,
the exact source syntax, the representation witness count, and the dependence
on every input coordinate. Preserve raw-input soundness for the proof checker.
In object-model arguments, preserve quantification over a model's entire
natural-number object rather than silently restricting to externally standard
numerals.

\subsection{Use several measures of success}

Record hand-maintained proof structure, generated source, declaration count,
source bytes, elaboration time, peak memory, and endpoint axioms separately.
A short line containing 36 rewrite lemmas can be more burdensome than several
well-named local facts. A three-line \code{decide} proof can be expensive to
replay. A new general theorem can be an architectural improvement even if its
first clients do not amortize the helper's source length.

\Needspace{7\baselineskip}
\section{Prioritized proposals}

\begin{longtable}{p{0.07\textwidth}p{0.35\textwidth}p{0.49\textwidth}}
\toprule
Order & Change & Why this target is concrete\\
\midrule
\endhead
1 & Replace private lookup inductions & Seven same-statement replacements,
five files, 40 lines removed; core replacements checked.\\
2 & Normalize half-radix digits algebraically & One quotient/remainder identity
replaces separate digit cases and exposes the parity argument.\\
3 & Share coefficient-row proof structure & 22 envelope lemmas contain 158
identical arithmetic leaf scripts; keep the coefficient data fixed.\\
4 & Respect earlier circuit register lemmas & Two rewrite lists shrink from
19 to 7 and 36 to 9 entries; generator and output change together.\\
5 & Add one assignment-transport law & Removes repeated \code{Sum}/\code{Fin}
coordinate extensionality in representation constructors.\\
6 & Share sum-of-squares and cost lemmas & Three identical private evaluator
lemmas and an already existing parameterized cost argument can be promoted.\\
7 & Share scanner simulation and invariants & Five scanner families repeat
iterator facts; the raw-input and decoded-list interfaces must stay distinct.\\
8 & Extract the common 99-node circuit tail & Exact source-data comparison
passes under a concrete renaming; one block proof can serve both families.\\
9 & Generalize DAG translation certificates & Reuse existing node-equation
uniqueness and reference relations for positive/natural conversions.\\
10 & Reuse object-model proof wrappers & A completeness/context wrapper occurs
26 times across 17 object-theory files.\\
11 & Use smaller views of existing window facts & The coherent \code{Facts}
record already exists; expose derived fields and named construction.\\
12 & Improve generated finite-check granularity & Reduce repeated unfolding
and make the trusted checking boundary explicit; benchmark before promotion.\\
\bottomrule
\end{longtable}

The ordering favors small, independently reviewable changes. The largest
architectural idea is number 8; it should follow the low-risk interface cleanup.
These are priorities, not estimated completion dates or claimed global savings.

\section{Immediate reuse: eliminate seven private lookup proofs}

The following files repeat elementary induction on a list or lookup index:
\src{AnnotatedRename}{, }
\src{AnnotatedBound}{, }
\src{AnnotatedProofChecker}{, }
\src{AnnotatedRuleSide}{, and }
\src{AnnotatedAssumption}{.}
Lean 4.32.0 already provides the required directions:

\begin{Lean}
theorem lookup_mem {α : Type u} {xs : List α}
    {n : Nat} {a : α} (h : xs[n]? = some a) : a ∈ xs :=
  List.mem_of_getElem? h

theorem exists_lookup {α : Type u} {xs : List α}
    {a : α} (h : a ∈ xs) : ∃ n, xs[n]? = some a :=
  List.getElem?_of_mem h
\end{Lean}

The supplied \code{lookup-cleanup.patch} replaces five proofs of the first
kind and two of the second. It leaves declarations, names, statements, and
call sites in place. The patch applies to the pinned source and removes
exactly 40 physical lines. \code{LookupCleanup.lean} checks the replacement
terms against the pinned core library; its axiom output is recorded.
This is a direct low-risk patch, although compiling the five modified
repository modules and their downstream endpoints is still a separate
integration check.\cite{lean-list}

This observation suggests a focused library-reuse pass. Search private
inductive lemmas about lists, finite sums, powers, and iteration before writing
new abstraction layers. Prefer replacing an established duplicate to creating
a parallel local theorem with slightly different notation. Do not generalize
every short private fact: a helper should have a clear reusable semantic role.

\section{Move renaming out of expression evaluation}

\src{Representation}{ already provides \code{Expr.eval_rename}.}
However, \code{Rep.reindex}, both sides of \code{Rep.conjunction}, and both
sides of \code{Rep.disjunction} repeat the same assignment comparison:
rewrite expression renaming, apply congruence, prove function equality, and
split a sum index.

The useful common law is
\[
 \operatorname{natAssignment}(v,w)\circ\operatorname{Sum.map}(f,g)
 =\operatorname{natAssignment}(v\circ f,w\circ g).
\]
It gives the semantic rule:

\begin{Lean}
theorem eval_rename_natAssignment
    (e : Expr (σ ⊕ α)) (f : σ → τ) (g : α → β)
    (v : τ → Nat) (w : β → Nat) :
    (e.rename (Sum.map f g)).eval (natAssignment v w) =
      e.eval (natAssignment (v ∘ f) (w ∘ g)) := by
  rw [Expr.eval_rename, natAssignment_comp_sumMap]
\end{Lean}

The assignment lemma itself is proved by \code{funext i; cases i <;> rfl}.
The left and right witness layouts, using \code{Fin.castAdd} and
\code{Fin.natAdd}, become instances. The supplied
\code{AssignmentTransport.lean} contains the general law and both instances.

Orient simplification toward removing a renaming from the expression evaluator
and putting it in the assignment. Do not install both directions as simp rules.
Keep these rules local or in a small named collection until their effect on
other evaluators is clear. Coordinate constructors should expose their semantic
interface before the larger arithmetic expression is normalized.

This preserves executable representation data. In particular, replacing a
\code{Rep} constructor with an existentially chosen representation would lose
the point of the current design: \code{Rep} carries a specific expression
and finite witness count in \code{Type}. The desired refactor changes proof
fields while preserving those supplied data.

\section{Replace half-radix case analysis with one identity}
\label{sec:radix}

\src{NativeCyclicRotation}{ defines a shifted digit by}
\[
 D_{\rho,s}(a,b)=(sa\bmod\rho)+\left\lfloor\frac{sb}{\rho}\right\rfloor,
\]
where $a=c(k)$ and $b=c(k-1)$. Its half-radix bounds analyze the possible
small digits separately. The arithmetic has a simpler normal form.

\begin{proposition}
For natural numbers $s,a,b$ with $s>0$,
\[
 D_{2s,s}(a,b)=s(a\bmod2)+\lfloor b/2\rfloor.
\]
Consequently, if $b\le3$, then $D_{2s,s}(a,b)\le s+1$.
If also $s\equiv0\pmod2$, then
\[
 D_{2s,s}(a,b)\bmod2=\lfloor b/2\rfloor.
\]
\end{proposition}

\begin{proof}
Multiplication by $s$ commutes with remainder modulo $2s$, and positive
multiplicative cancellation gives $\lfloor sb/(2s)\rfloor=\lfloor b/2\rfloor$.
Now $a\bmod2\le1$ and $b\le3$ imply $\lfloor b/2\rfloor\le1$.
When $s$ is even, the first summand has even parity and the second is less
than 2, giving the final equality.
\end{proof}

The exact Lean proof of the identity is short:

\begin{Lean}
theorem step_eq (s a b : Nat) (hs : 0 < s) :
    step s a b = s * (a % 2) + b / 2 := by
  unfold step
  rw [Nat.mul_comm 2 s, Nat.mul_mod_mul_left,
    Nat.mul_div_mul_left _ _ hs]
\end{Lean}

\code{HalfRadix.lean} checks the identity, bound, and parity statement using
only Lean core and its proof-producing \code{omega} tactic. This is stronger
than the existing local bound: the current digit $a$ needs no upper bound;
only the predecessor $b$ is constrained. Existing public statements can remain
as corollaries, so callers need not change.

The distinction between generality and integration matters. This file proves
the scalar arithmetic used by the repository definition. It does not replay
the entire cyclic-rotation dependency chain or the later phase argument.
Specialization to $s=2^{88}$ is routine, but should still be compiled in the
original module before committing the replacement.

\section{Compress the coefficient-row proofs without hiding their data}
\label{sec:rows}

\src{NativeCoefficientBounds}{ contains 22 private
\code{rowEnvelope_0}, \ldots, \code{rowEnvelope_21} lemmas.}
Between them, 158 leaves execute the same pattern:

\begin{Lean}
rw [abs_le]
constructor <;>
  linarith only [h0, h1, h2, h3, hsq, hmf, hkp]
\end{Lean}

Each row selects a finite piecewise coefficient expression, all bounded by
the same envelope. The large source block repeatedly copies the row function,
introduces the same three nonnegativity facts, splits its conditional tree,
and runs the same linear reasoning. The row data are fixed independently in
\src{NativeQuadraticParent}{.}

The direct refactoring is one theorem with a row index $k:\operatorname{Fin}(22)$.
Case-split $k$, reduce only the selected coefficient row, decompose the conditional
bounds, and close each arithmetic leaf using the shared facts. The
pair index stays symbolic; there is no need to enumerate all 528 unordered
pairs simply to prove a uniform magnitude bound.

\begin{Lean}
-- Schematic proof shape; the candidate is not kernel-verified.
have hsq := sq_nonneg (ports 0)
have hmf := mul_nonneg h3 h0
have hkp := mul_nonneg h1 h0
fin_cases k <;>
  dsimp only [rowCoefficient, Fin.reduceFinMk, Matrix.cons_val] <;>
  repeat' apply abs_ite_le
all_goals
  rw [abs_le]
  constructor <;>
    linarith only [h0, h1, h2, h3, hsq, hmf, hkp]
\end{Lean}

Controlled unfolding is essential. A global \code{simp} over the parent circuit
could unfold much more than these finite coefficient rows. An initial
\code{split_ifs} approach exceeded the simplifier's step limit on this input.
The revised proof keeps the existing small \code{abs_ite_le} lemma and descends
through the conditional bound structurally, without adding branch hypotheses.
Vector indexing is reduced before this step. The final verification receipt
distinguishes checking against exact extracted definitions from a full-module
replay.

The supplied whole-file candidate changes 679 physical lines to 130, a
549-line source reduction. This figure is a candidate source measurement;
its verification scope is stated separately. Removing private row wrappers
is permitted only after confirming they are not used externally. Keep the
public \code{abs_rowCoefficient_le_envelope} and later envelope theorem interfaces.

\subsection{Do the same for coefficient naturality}

\src{NativeNinePolynomial52}{ contains 22 private
\code{mapRowCoefficient} specializations.}
They establish that a ring homomorphism commutes with each row's coefficient
recipe. The mathematical reason is uniform: coefficients use integer literals,
addition, subtraction, multiplication, powers, and conditionals whose tests
depend on indices rather than ring values.

One finite-row dispatcher with the homomorphism simp laws can replace the
duplicated literal functions. A longer-term design stores the coefficient
recipe as syntax and derives naturality once by structural induction. Before
introducing new syntax, use the existing expression evaluators where suitable.
The independent theorem tying those recipes to the actual source rows must
remain: redefining the source row to equal its encoded coefficient sum would
remove the verification obligation instead of proving it.

The accompanying mapping candidate reduces its full file from 456 to 153
physical lines, a 303-line deletion. The combined isolated coefficient harness
was not verified: its last capped run timed out without diagnostics. Both
reductions therefore remain candidate source measurements, not completed
kernel-checked migrations.

\subsection{A stronger abstraction: a finite majorant certificate}

The envelope has the nonnegative monomial basis
\[
 m=(\beta^2,\ M_F\beta,\ M_C,\ M_F,\ \beta,\ K\beta,\ K,\ 1)
\]
with weight vector $b=(2,2,1,1,2,1,1,8)$. Each coefficient branch is an integer
linear combination $\sum_i a_i m_i$ with small integer $a_i$. Prove once that
$m_i\ge0$ and $|a_i|\le b_i$ imply
\[
 \left|\sum_i a_i m_i\right|\le\sum_i |a_i|m_i
 \le\sum_i b_i m_i.
\]
Then each concrete branch needs only a finite small-integer vector certificate.
This can serve multiple coefficient packages and reduces repeated arithmetic
reasoning more deeply than a compact tactic script. It is a proposed design,
not an implemented result. The basis must retain the existing linear dependence
on the large transport port $K$; an easier bound quadratic in $K$ would not
preserve the downstream published scale.

\section{Use established register facts as proof boundaries}
\label{sec:cuts}

\src{SplitCircuitRegisters59}{ unnecessarily expands some registers whose
semantic values have already been proved.}
The corresponding Gap-family file demonstrates the better discipline. Two
particularly clear targets are:

\begin{center}
\begin{tabular}{lrrrr}
\toprule
Theorem & Rules before & Rules after & Bytes before & Bytes after\\
\midrule
\code{natural_sourceSRegister} & 19 & 7 & 744 & 487\\
\code{natural_sourceRankRegister} & 36 & 9 & 1,142 & 555\\
\bottomrule
\end{tabular}
\end{center}

The first should use \code{natural_sourceHighRegister}. The second should
reuse the previously established source-$S$, source-$T$, and block-length
identities. The generated Lean file shrinks by 844 UTF-8 bytes, from 14,604
to 13,760. Both theorem declarations remain eight physical lines because the
rewrite lists were printed on long lines. This is a reduction in exposed
dependencies, not a measured physical-line saving.

The supplied patch changes the generator and generated output together.
As theorem entries are scheduled, it snapshots the currently available
semantic cuts and adds each new cut only for later entries:

\begin{Lean}
entries.append((name, idx, expr, dict(cuts), False, defs))
cuts[idx] = name
\end{Lean}

Copying the dictionary is necessary. If an entry retained a reference to the
same mutable dictionary, its generated proof could see a later theorem or
its own theorem. The improvement depends on respecting source order, not only
on finding shorter rewrite lists. Regeneration, Lean checking, and downstream
checking are three distinct checks; the archive labels them separately.

\section{Extract the shared 99-node free-carrier tail}
\label{sec:tail}

The most specific larger opportunity is visible in both the mathematical
specification and the literal circuit data. The \code{ratioValue} definitions
in \src{SplitNinePolynomial59}{ and }
\src{GapNinePolynomial66}{ are byte-identical: 24 lines, 612 bytes.}
More strongly, the complete natural-circuit suffixes at Split indices
$416,\ldots,514$ and Gap indices $385,\ldots,483$ are identical after reference
renaming. This is 99 nodes: 60 multiplications, 15 subtractions, 17 additions,
and 7 literal nodes. There are no input or fixed-coefficient leaf nodes inside
the tail.

Internal references map by $j\mapsto j-31$. The boundary references map as
follows:

\begin{center}
\begin{tabular}{rrl}
\toprule
Split register & Gap register & Role\\
\midrule
7 & 7 & Internal input $j$\\
9 & 9 & One\\
10--17 & 10--17 & Shifted external inputs\\
21 & 87 & Two\\
409 & 378 & Threshold $N$\\
414 & 383 & Rank\\
415 & 384 & Margin scale\\
\bottomrule
\end{tabular}
\end{center}

The archive's \code{inspect_shared_tail.py} compares the pinned generator
node dictionaries exactly and emits a JSON receipt. That comparison identifies
the target; it is not yet a Lean proof that a new block abstraction preserves
evaluation. The script's source inputs and commit are explicit.

\subsection{A suitable common interface}

Define a small straight-line block whose operand references either a boundary
port or an earlier local node. Its 14 boundary ports include the constants one
and two; its semantic specification is the common free-carrier finalizer.
Then establish four facts:

\begin{enumerate}
\item A block evaluator agrees with insertion into a parent topological DAG
when the boundary values agree.
\item This one 99-node block evaluates to the common \code{ratioValue}.
\item Each existing parent suffix is the exact reference-instantiation of
the block.
\item The existing front-end register lemmas supply the boundary values.
\end{enumerate}

Insertion should instantiate references directly, adding no copy gates.
Retain equality with each pinned parent node list, the operation counts, and
the output index. Source-level sharing must not change the measured equation
or hide a changed literal coefficient tuple.

The typing issue is real: Split has 12 fixed inputs and Gap has 11, while the
existing \code{NodeRel} is formulated for equal input and fixed-input arities.
Since the common tail does not read fixed-input leaves, a boundary-port block
captures the exact shared object without coercing the two parent node types.

Put the common block below both specializations in the import graph. It should
depend on generic DAG semantics and common finalizer algebra, not on either
entire specialized circuit. A module cycle would cancel much of the intended
simplification.

This proposal reduces duplicated proof about 99 nodes. It does not remove
99 arithmetic operations from either Diophantine construction, and no such
size improvement is claimed. The new generic machinery must be charged against
any eventual source reduction.

\section{Share scanner proofs at the correct semantic boundary}
\label{sec:scanners}

Five families repeat the same iteration facts:
\src{AnnotatedRenameScan}{, }
\src{AnnotatedBoundScan}{, }
\src{AnnotatedLineScan}{, }
\src{AnnotatedContext}{, and }
\src{AnnotatedAxiomContext}{.}
Their 15 proofs of \code{scan_zero}, \code{scan_succ}, and \code{scan_done}
contain 72 proof-body lines. Mathlib already has the corresponding
\code{Function.iterate} facts. A one-time bridge from the current
\code{Nat.rec} convention can reuse them.\cite{mathlib-iterate}

The larger repetition is the simulation of an encoded list by a numeric
scanner. The following interface isolates it without replacing the supplied
numeric program:

\begin{Lean}
structure Encoding (next : A → H → Option H) (step : S → S) where
  state : List A → H → S
  reject : S
  reject_fixed : step reject = reject
  done_fixed : ∀ h, step (state [] h) = state [] h
  cons_step : ∀ row rows h,
    step (state (row :: rows) h) =
      match next row h with
      | none => reject
      | some h' => state rows h'
\end{Lean}

For the associated typed list executor \code{run}, the common theorem is:
\[
 |\mathit{rows}|\le f\quad\Longrightarrow\quad
 \operatorname{iterate}(\mathit{step},f,
     \mathit{state}(\mathit{rows},h))
 =\operatorname{result}(\operatorname{run}(\mathit{rows},h)).
\]
The proof is induction on the list. At an empty list, the state is fixed;
at a rejected row, rejection is fixed; at a successful row, consume one unit
of fuel and apply the induction hypothesis. The prelude-only
\code{ScannerCore.lean} checks this theorem and the accompanying invariant
lemmas without additional axioms.

\subsection{Do not weaken arbitrary-code soundness}

The existing \code{AnnotatedLineScan.scan_sound} applies to arbitrary natural
codes and arbitrary fuel. A theorem about \code{state rows history} covers
only states arising from the chosen encoding. The latter cannot replace the
former. Keep an invariant on raw numeric states, prove one-step preservation
using the actual local checker, and lift it through the common iteration
theorem. Use encoded-list simulation for completeness and sufficient-fuel
arguments. This separation removes repeated induction without introducing a
parser assumption.

The same distinction applies to finished states. The generic interface does
not assume that a rejecting state and every completed state are distinct.
Any distinction needed by the final acceptance test must remain an explicit
property of the concrete numeric implementation.

\subsection{Share typed execution and proof fragments}

\src{AnnotatedRename}{, }
\src{AnnotatedBound}{, and }
\src{AnnotatedProofChecker}{ repeat list execution, soundness, and production
composition.}
The checked generic \code{run_append} and \code{run_invariant} cover the
first layer. A small history-specific layer can cover \code{Produces.append}
and \code{Produces.step}. Preserve the reversed history convention: executing
$p$ and then $q$ builds history in the order $h_q\mathbin{++}h_p$ because parent
references count backwards. A superficially cleaner forward-list composition
would change the checker semantics.

\subsection{Factor the supplied PR iterator, not just its existence proof}

The bound, line, and axiom-context scanners also repeat program construction
from a base, a step, and primitive recursion. A parameterized
\code{iterateProgram} could expose base-history initialization and fuel
projection as inputs. This affects executable program data. Preserve or prove
the exact \code{.val} equality, arities, projections, and node-count formulas.
Replacing the constructor by an existential proof of primitive recursiveness
would discard information required later by the Diophantine representation.

\section{Share residual-system facts before merging compilers}
\label{sec:squares}

Three files contain the same private \code{eval_sumExpressions} induction:
\src{DAGQuartic}{, }
\src{NaturalDAGQuartic}{, and }
\src{QuadraticRegisterSystem}{.}
Publish it in the syntax layer and prove once, for an arbitrary finite
residual family $r_i$, that
\[
 \left(\sum_i r_i^2\right)(v)=0
 \quad\Longleftrightarrow\quad
 \forall i,\ r_i(v)=0
 \qquad\text{over }\mathbb Z.
\]
Each compiler's \code{expression_zero_iff} then only unfolds its residual
family and invokes this interface. The concrete proof is in the proposed
\code{CircuitCoreRefactors.lean}; its checked status is listed separately.

The codomain restriction is substantive. Over a general commutative ring,
a sum of squares need not vanish only when each square does. For example,
$1^2+i^2=0$ over the complex numbers. The internal-omega theorem in
\code{ObjectResidualSquares} uses the object theory's positive/negative
semiring interpretation and order. It cannot be replaced by the integer
lemma without an appropriate proof of transfer.

\subsection{Promote the parameterized cost bound already present}

\src{QuarticCost}{ specializes a sum-expression operation bound to a fixed
per-expression budget.}
\src{DAGQuartic}{ already proves the parameterized version privately.}
Promote the latter and retain the old result as a specialization:
\[
 \bigl(\forall e\in E,\ \operatorname{binaryOps}(e)\le b\bigr)
 \Longrightarrow
 \operatorname{binaryOps}(\operatorname{sumExpressions}(E))
 \le (b+1)|E|.
\]
This is reuse of an existing argument rather than new mathematics.
Keep the sharper special treatment of the final output residual: applying
a coarse uniform bound everywhere can worsen the stated operation bound.

\subsection{Keep signed and nonnegative wire realizations distinct}

The signed and nonnegative DAG compilers share gate semantics and degree
lemmas, but realize wires differently. Signed values use positive and negative
natural parts; the single-wire natural compiler requires a nonnegativity
invariant. Common RHS syntax and residual lemmas are good targets. A universal
compiler interface should carry the required wire realization theorem rather
than silently adding nonnegativity to the signed theorem. Exact witness counts
remain separate corollaries of the two encodings.

\section{Use a common DAG invariant theorem selectively}

\src{ArithmeticDAGSemantics}{ already proves the local node equation and
uniqueness of a value table satisfying those equations.}
\src{ArithmeticDAGDegree}{ and }
\src{NaturalDAGQuartic}{ repeat the corresponding topological strong induction
for degree and nonnegativity.}

A general invariant has the following mathematical form. Let $P(i,r)$ be a
property of register $i$ and value $r$. If each gate at $i$ preserves $P$
whenever $P$ holds at all its dependency registers, then $P$ holds at every
register in the actual run. Topology supplies the strict predecessor bound;
\code{run_node_equation} supplies the gate equation.

\begin{Lean}
-- Essential local premise of run_invariant:
∀ i < nodes.length, ∀ values,
  (∀ j ∈ (nodes.getD i (.literal 0)).dependencies,
    P j (values j)) →
  P i ((nodes.getD i (.literal 0)).evalAt p x values)
\end{Lean}

The supplied light \code{DAGInvariant.lean} and larger
\code{CircuitCoreRefactors.lean} show proposed clients for degree bounds,
semantic-table uniqueness, and nonnegative runs. This also supports future
range bounds and symbolic normal-form certificates.

This is not an immediate net line-saving result in the measured examples.
Three original declarations occupy 42 lines and 2,083 bytes; a proposed helper
plus three replacements occupy 45 lines and 2,296 bytes under the supplied
formatting. Longer example names contribute to the latter figure. The benefit
is concentrating the induction invariant and topology bookkeeping. Adopt it
when additional clients justify the common interface, or when it makes a
particularly fragile proof substantially easier to maintain.

\section{Generalize the existing translation and reflection interfaces}

\subsection{Positive/natural translation is a simulation, not a ring map}

\src{SplitCircuitTranslation59}{ and }
\src{GapCircuitTranslation66}{ follow the same proof pattern: handle a short
initial prefix, handle one operand-swapped multiplication, then transport all
other node equations through a reference map.}
They use an exceptional register 7 and an input shift. The special swapped
nodes occur at Split 452 and Gap 421.

The repository already has \code{ArithmeticDAG.run_map} and
\code{value_map}. These are valuable for genuine ring homomorphisms, but
$x\mapsto x-1$ is not a ring homomorphism, and these two translations change
node lists. Do not try to solve this problem by applying naturality with that
shift.

Instead, parameterize the existing simulation proof by its reference map,
initial-prefix facts, exceptional-value facts, and node-relation certificate.
Extend \code{NodeRel}, if useful, to allow operand swapping for addition and
multiplication only. Then use existing \code{run_eq_of_node_equations} or
\code{value_eq_of_node_equations} for the conclusion. Each specialization
retains small closed finite certificates and its boundary equations.

The name ``natural'' describes the external input domain in these files;
the node lists still contain negative literals and subtraction. They therefore
do not satisfy \code{NaturalDAGQuartic.Allowed} merely because of their names.

\subsection{Do not replace efficient local lemmas by uncontrolled unfolding}

The step files contain 999 ordinary step theorems and 193 power theorems.
An indexed theorem or a proof-producing local command could generate only
needed specializations. However, the current named lemmas provide cached
rewrite boundaries. Removing them can force repeated reductions of large
literal lists. The meaningful experiment measures cold and incremental replay
of actual clients, not only the emitted file size.

\src{ArithmeticDAGPowers}{ already provides power certificates, reference
relations, and rewrite soundness.}
\src{SplitJointPower59}{ already uses that architecture.}
Extend those interfaces to blocks and local identities rather than introducing
an unrelated reflection framework with a second semantics.

\subsection{A small sparse-row normalizer is a better first reflection target}

The row-identity families in Gap, Split, and Native each compare low-degree
source residuals with coefficient representations. A normalized sparse
quadratic syntax is a bounded, reviewable target for a certificate checker.
Check its equality to the independently supplied source recipe. Avoid expanding
the full high-degree finalizer into a polynomial normal form: that destroys
the sharing which the circuit representation was designed to preserve.

\subsection{Preserve cancellation-aware degree arguments}

Numerical degree labels give upper bounds. Algebraic rewriting can give a
strictly better bound by exposing cancellation. The existing
\src{FreeCarrierDegree}{ and }
\src{GapFinalizerDegree66}{ deliberately exploit such structure.}
A shorter certificate for a weaker naive degree bound is not a replacement
for the stronger theorem. A promising block checker first proves the rewrite
identity and then applies the existing weighted-degree interpretation to that
chosen normal form. Any Pell or other algebraic hypotheses remain explicit.

\section{Represent elementary row operations once}

\src{SplitCertificate59}{ proves \code{transformed_zero_iff} by explicitly
handling all 22 coordinates in both directions.}
Most coordinates are unchanged; two rows are modified by adding a different
row. The proof is an instance of a very small general fact.

For a row family $r:I\to R$ and distinct indices $d,s$, define
\[
 (\operatorname{addAt}_{d,s}r)(k)=
 \begin{cases}r(d)+r(s)&k=d,\\r(k)&k\ne d.\end{cases}
\]
Then
\[
 (\forall k,\ (\operatorname{addAt}_{d,s}r)(k)=0)
 \quad\Longleftrightarrow\quad (\forall k,\ r(k)=0).
\]
The forward direction first reads the unchanged source row $s$, obtaining
$r(s)=0$, and then reads the destination row. No cancellation or inverse
operation is necessary: a type with addition and the usual zero laws suffices.
Distinctness of $d$ and $s$ is essential to this argument.

\code{RowOperations.lean} checks this generic theorem and the corresponding
fact for a surjective reindexing. The intended concrete refactoring expresses
the transformed family as two elementary row additions, then invokes this
theorem twice. The final reordering uses the existing
\code{rowOrder_surjective}. The equality between the actual transformed
syntax and those elementary operations remains to be established; it is the
finite binding certificate, not an assumption to omit.

This is preferable to introducing a full matrix library solely for two row
additions. If later transformations include arbitrary invertible linear
combinations, a matrix-based zero-set preservation theorem becomes appropriate.
The current case supports a much smaller interface.

\section{Simplify object-model proof transport without changing strength}
\label{sec:object}

The exact expression \code{completeness_inf_context ZFCax_s []} appears
26 times across 17 \code{Object*.lean} files. The repeated proof shell obtains
a model, extracts its ZF axioms, establishes validity of the sealed formula,
and invokes completeness. Good targets include
\src{ObjectNaturalExpression}{, }
\src{ObjectCodeClosure}{, }
\src{ObjectCodeBound}{, }
\src{ObjectPairingGraph}{, and }
\src{ObjectResidualSquares}{.}

Prove a wrapper taking the same model-universe quantification and semantic
validity hypothesis and returning the same sealed \code{BProv} conclusion.
Then every caller can focus on its specific arithmetic or coding argument.
This theorem should be derived from the current completeness and sealing
lemmas; it adds no axiom and gives no license to transfer arbitrary external
theorems about standard naturals into ZFC.

\subsection{An especially small arity-polymorphism improvement}

In \src{ObjectQuarticCompiler}{, \code{zfc_proves_compiler} proves two versions
of the same environment-invariance fact.}
Its local helper is tied to the input arity $m$, so the compiled expression's
different arity needs a separate \code{closedTarget} proof. Generalize

\begin{Lean}
have closed (f : Expr (Fin m)) : ... := ...
\end{Lean}

to

\begin{Lean}
have closed {n : Nat} (f : Expr (Fin n)) : ... := ...
\end{Lean}

and use it at both arities. The body already uses the arity-independent
\code{sat_noZeroForm_iff} and \code{Sat_sentence_inv}. This proposed edit is
source-inspected; the package does not claim a fresh build of its full
object-theory dependency chain.

\subsection{Build on the algebraic interfaces already present}

The project already has \code{ObjectOmegaSemiring},
\code{ObjectNaturalOrder}, \code{GatesSemiring}, and \code{QuarticSemiring}.
The next step is to share model-independent finite-index and renaming facts
between their callers. For example, conjunction of row families has the same
index argument in \code{ObjectQuadraticSystem} and
\code{QuadraticRegisterComposition}.

Two constraints affect the design. First, directly importing
\code{ObjectNaturalOrderedSemiring} into \code{ObjectResidualSquares} would
create a cycle through \code{ObjectNaturalDiscrete}. Extract lower-level
algebraic lemmas with only the needed assumptions or deliberately reorder the
modules. Second, \code{ObjectNaturalExpression.evalRel_unique} currently does
not assume every input lies in omega. A replacement proof that requires all
inputs to be natural would weaken that theorem. Totality and uniqueness need
not use exactly the same abstraction.

\section{Improve the current-window API where coherence already exists}

\src{CurrentPolynomialWindow2011}{ already contains a substantial
\code{Facts} structure tying one source assignment, one window length,
one Boolean word, one carry function, and one digit function together.}
That design should be reused. The proposed improvement is to expose smaller
proved views and derived quantities, not to introduce a competing record.

Examples include a named temporal-shift expression, a view containing the
four carry/mask facts needed by a local decoding lemma, and named fields for
the source-coordinate transport. Construct \code{Facts} using named fields
instead of a long positional tuple when this helps future edits. The latter
does not necessarily shorten source, but makes an added field easier to audit.

\src{CurrentDecodedWindow2011}{ correctly derives count and phase from the
same arithmetic facts before invoking the cell decoder.}
Keep that order. A ``simplified'' cell interface that assumes occupation or
phase before the relevant polynomial theorem supplies them would make the
composition circular. In particular, the decoded-cell alternatives still allow
empty cells; the new interface must not quietly assume that every cell selects
a tile.

At the arithmetic boundary, derive reusable conversion facts once. Natural
subtraction is truncated, integer subtraction is not, and \code{ZMod} values
are representatives modulo the full window length. An unconditional rewrite
between these interpretations is usually wrong. Good helper statements carry
the exact positivity or range premise that makes the conversion valid.
Then use \code{omega} for the remaining linear natural-number arithmetic and
\code{ring} for equalities after entering the ring domain.

This keeps representation changes separate from mathematical content. It is
also more robust than adding long lists of low-level cast and subtraction
lemmas to every downstream \code{simp} call.

One especially local example is congruence weakening. The result
\code{CurrentClause2011.reflected_weight_mod} works modulo 2048;
\code{CurrentEarlyPadding2011.reflected_mod_sixteen} repeats the same
selector/nonselector analysis modulo 16. The latter should follow through
\code{Nat.ModEq.of_dvd}, since $16\mid2048$. The early-padding module already
imports the clause module. The two local selector functions have the same
defining equations; their definitional identification or a shared definition
must be checked. This is a source-inspected one-line-wrapper opportunity,
not a replayed project edit.

\section{Finite checking and generated proof design}

The desired trust boundary is ordinary Lean proof checking. A generator may
produce a finite table, a reference map, or a proposed algebraic certificate;
Lean must verify the assertion that table is used to justify. Keep certificate
data independent of the conclusion and preserve the current endpoint axiom
audits.

Several practical experiments are worthwhile:

\begin{enumerate}
\item \textbf{Split large closed checks into stable pieces.} Compare a single
\code{decide} over a large nested finite predicate with a small indexed family
of checks or a partitioned certificate. This may change resource use, but
neither short source nor many small cases guarantees faster replay.
\item \textbf{Reduce data selection before symbolic branching.} For row
coefficient proofs, choose the finite row and simplify vector indexing before
splitting conditionals. Otherwise a tactic may inspect branches belonging to
unselected rows.
\item \textbf{Keep high-level cuts explicit.} Do not make large register
expansion lemmas global simp rules. Use local lists of established semantic
facts and prove small ring identities at the resulting boundaries.
\item \textbf{Edit the generator with its output.} Check byte-for-byte
regeneration and compile the emitted proof. A passing generator test alone
does not certify its mathematical output.
\item \textbf{Retain compatibility wrappers when useful.} Public theorem
names referenced by checks can remain as thin corollaries of a generic proof.
Private wrappers can be removed after confirming their scope.
\end{enumerate}

These recommendations do not require replacing kernel reduction with
\code{native_decide}. The current source avoids that tactic. A tactic or
certificate generator should emit ordinary checked proof terms if the existing
axiom footprint is to be retained. Changes to trust assumptions should be
explicit design decisions, not incidental source-shortening techniques.

\section{Verification results and limitations}
\label{sec:validation}

The archive separates three kinds of evidence: kernel checks of new general
lemmas, exact source-data checks, and source-inspected proposals. A success in
one category is not promoted to another.

\begin{longtable}{p{0.35\textwidth}p{0.57\textwidth}}
\toprule
Artifact & Verified scope\\
\midrule
\endhead
\code{ScannerCore.lean} and audit & Seven general scanner/iteration theorems
compiled with Lean 4.32.0; all seven have empty axiom lists. Concrete scanner
adapters are not yet integrated.\\
\code{LookupCleanup.lean} & Both exact polymorphic replacement statements
compiled; each uses only \code{propext}. Seven-edit patch applies to the pinned
source; the five full modified modules were not rebuilt.\\
\code{HalfRadix.lean} & Identity, bound, and parity theorems compiled. Axiom
lists use only \code{propext} and, for the bound/parity proofs, \code{Quot.sound}.\\
\code{AssignmentTransportCore.lean} & Polymorphic \code{Sum.elim}/\code{Sum.map}
composition law compiled; its axiom list is \code{Quot.sound}. Repository-specific
clients are separate candidates.\\
\code{RowOperations.lean} & Elementary row-addition and surjective-reindexing
theorems compiled. The concrete Split certificate binding is still proposed.\\
\code{inspect_shared_tail.py} & Pinned JSON inputs and regeneration of both
circuit/step families checked; 99-node reference-renamed equality passes.
No common-block Lean semantic theorem is claimed.\\
\code{split-register-cuts.patch} & Generator/output source transformation and
unchanged theorem statements checked. Full generated-module replay remains
required.\\
Coefficient-row candidates & Last combined isolated-definition replay timed
out after 90 seconds without diagnostics. The 549- and 303-line deletions are
unverified candidates. Exact scope is recorded in
\code{evidence/arithmetic/verification.json}.\\
Other circuit/model prototypes & Source-grounded proposals; their file-specific
status is recorded in \code{evidence/verification.json}. No whole-project
refactoring or build is claimed.\\
\bottomrule
\end{longtable}

\subsection{Resource-limited checks are not performance measurements}

The pinned compiler and dependency cache were obtained, and several upstream
project modules were built. An attempted targeted build later exited with
code 137 while building \code{SelectiveCertificate52}. The environment's shared
memory limit was 8 GiB and was under concurrent load. This is a resource-limited
replay, not evidence that the original proof is mathematically wrong or that
this particular certificate intrinsically requires more than 8 GiB.

A repository-dependent assignment-transport example also aborted under an
explicit 1 GiB Lean cap before any theorem diagnostic. An import-only probe
for the arithmetic harness exceeded that cap as well. These observations do
not establish a relative cost between the proposed and original proofs.
No timing speedup or memory reduction is claimed in this report.

\subsection{What the delivered evidence does establish}

The elementary general lemmas are actual checked Lean code. The source census
and fixed-data comparisons are reproducible. The patches and prototypes give
specific theorem-level changes that can be reviewed in the original project.
The package also identifies where full imported-module replay, concrete adapter
proofs, and endpoint audits remain necessary before integration.

\section{A concrete migration plan}

\textbf{First change set: library reuse and local algebra.}
Apply the lookup patch and run the affected checks. Add the half-radix identity
and retain the original bound/parity names as corollaries. Add the assignment
transport rule at the existing representation layer. Generalize the local
arity-polymorphic helper in \code{ObjectQuarticCompiler}. These changes have
small semantic surfaces and can be reviewed separately.

\textbf{Second change set: repeated finite proofs.}
Validate the row-envelope replacement against the original imports; preserve
the exact coefficient definition and the public theorem statements. Replace
the coefficient-mapping family with one naturality proof. Bind the actual
Split row transformation to the small row-operation interface. Promote the
existing shared sum-of-squares and cost lemmas. Compile downstream public
theorems after each step.

\textbf{Third change set: shared scanner and circuit infrastructure.}
Integrate the scanner interface first in RenameScan and BoundScan, where the
common shape is clearest, then LineScan and AxiomContext. Keep all-code
soundness and typed completeness separate. For the circuit family, apply the
semantic-cut generator change and then extract the 99-node tail with exact
parent-list binding theorems. Generalize positive/natural translation after
the block reference interface has settled.

At each step compare exact signatures, executable program data, witness counts,
and endpoint axioms. For generated code, run its reproducibility check. Measure
resource use on the same machine under comparable load only after the relevant
modules replay successfully. Report net size after counting new helper modules,
adapters, and compatibility wrappers.

\appendix
\section{Archive guide and reproduction}

The principal files are \code{report.tex} and \code{report.pdf}.
\code{README.txt} gives the file inventory and runnable commands.
\code{lean/checked/} contains checked small examples; \code{lean/candidates/}
contains proposed repository-dependent integrations. Patches, source-data
checkers, and evidence are in \code{patches/}, \code{scripts/}, and
\code{evidence/}. Detailed family notes are in \code{notes/}.

The entire report is one LaTeX source file. Build it with XeLaTeX twice:

\begin{Lean}
xelatex -interaction=nonstopmode -halt-on-error report.tex
xelatex -interaction=nonstopmode -halt-on-error report.tex
\end{Lean}

Use the repository's pinned toolchain and manifest for code replay. The exact
Lean release is \code{leanprover/lean4:v4.32.0}; its binary reports commit
\code{8c9756b28d64dab099da31a4c09229a9e6a2ef35}.
The supplied driver checks the standalone core files in a temporary output
directory and records command results. Repository-dependent candidates require
their imported modules to have been built; they are not advertised as
independently self-contained projects.

\begin{Lean}
python3 scripts/audit_sources.py /path/to/ProveIt /tmp/source-audit
python3 scripts/inspect_shared_tail.py /path/to/ProveIt
python3 scripts/check_core.py --lean /path/to/lean --out /tmp/core-checks
git -C /path/to/ProveIt apply --check /path/to/lookup-cleanup.patch
\end{Lean}

Inline code snippets omit surrounding namespaces and ambient variables when
that improves readability. The complete artifact files and their recorded
compiler results determine checked status, not the typography of a snippet.
The archive manifest records hashes of deliverables. No binary compiler,
Mathlib cache, or unrelated repository material is bundled.

\begin{thebibliography}{9}
\bibitem{readme}
ProveIt contributors.
\href{https://github.com/VladimirReshetnikov/ProveIt/blob/60e9c5d620c6c387d932876cfa2d9376f64ea56c/Computability/DiophantineConsistency/README.md}{DiophantineConsistency README}.
Pinned snapshot \code{60e9c5d620c6}, 26 October 2026.

\bibitem{plan}
ProveIt contributors.
\href{https://github.com/VladimirReshetnikov/ProveIt/blob/60e9c5d620c6c387d932876cfa2d9376f64ea56c/Computability/DiophantineConsistency/LEAN_EQUIVALENCE_PLAN.md}{Lean equivalence plan}.
Same snapshot. Describes the exact checked endpoints and outstanding composition
obligations.

\bibitem{repo}
ProveIt contributors.
\href{https://github.com/VladimirReshetnikov/ProveIt/tree/60e9c5d620c6c387d932876cfa2d9376f64ea56c/Computability/DiophantineConsistency/Lean}{Active Lean source tree}.
Specific modules are linked throughout this report; reproducible file hashes,
source spans, and imports accompany the archive.

\bibitem{lean-list}
Lean contributors.
\href{https://github.com/leanprover/lean4/blob/8c9756b28d64dab099da31a4c09229a9e6a2ef35/src/Init/Data/List/Lemmas.lean}{Core list lemmas},
Lean 4.32.0. In particular \code{List.mem_of_getElem?} and
\code{List.getElem?_of_mem}.

\bibitem{mathlib-iterate}
Mathlib contributors.
\href{https://github.com/leanprover-community/mathlib4/blob/81a5d257c8e410db227a6665ed08f64fea08e997/Mathlib/Logic/Function/Iterate.lean}{Function iteration lemmas}.
Exact Mathlib revision from the repository manifest.

\bibitem{manifest}
ProveIt contributors.
\href{https://github.com/VladimirReshetnikov/ProveIt/blob/60e9c5d620c6c387d932876cfa2d9376f64ea56c/lake-manifest.json}{Dependency manifest}
and
\href{https://github.com/VladimirReshetnikov/ProveIt/blob/60e9c5d620c6c387d932876cfa2d9376f64ea56c/lean-toolchain}{toolchain file}.
These pins determine the environment used for the reported checks.
\end{thebibliography}
\end{document}


